% paper/sections/08_2_fccd_bf.tex
%
% §8.4  FCCD 面ジェットによる BF sub-system 統合
%       （H-01 歴史的診断は付録 §B.4）
%
% ═══════════════════════════════════════════════════════════════════════════════
\subsection{FCCD 面ジェットによる BF sub-system 実装}
\label{sec:fccd_bf_sub_system}
% ═══════════════════════════════════════════════════════════════════════════════

節 \ref{sec:fccd_def} で導出した FCCD 面ジェット
$\FaceJet{u}=(u_f,u'_f,u''_f)$ は，BF sub-system の
P-1..P-4 を\textbf{同時に}満たす最小単位であり，
4 役割 (R1)-(R4) の (R3) 分相 PPE 楕円演算子への適用例として位置付けられる．
（旧設計（節点 CCD + 事後面平均）で生じていた $\Ord{\Delta x^2}\cdot|d\log J/dx|$
残差は歴史的に H-01 として診断された；付録~\ref{sec:fccd_motivation_h01} を参照．）

% -------------------------------------------------------------------------------
% H-01 診断の一次定義は §\ref{sec:fccd_motivation_h01}（§4.5）にある．
% 本節では FCCD 面ジェットによる救済に直ちに進むため，診断再掲は外部参照に委ねる．
% 旧 \subsubsection{H-01 診断の再掲と FCCD 救済} は
% paper/archive/v1_pre_sp_core/snippets/h01_recap_08_2.tex に退避．
% -------------------------------------------------------------------------------
\phantomsection
\label{sec:h01_diagnosis_fccd_remedy}% 外部参照保持：10_5/08c/10_3/appendix_gpu_fvm

本稿の構成では {\FCCD} 面ジェット $\FaceJet{u}=(u_f,u'_f,u''_f)$ が
圧力勾配・界面力・HFE 上流点・GFM ジャンプ条件の 4 経路で共有される（(R1)-(R4)）．
これにより面フットプリントは 4 経路で一致し，
\begin{equation}
  \mathcal{G}p := D^\FCCD p,
  \qquad \bnabla\psi := D^\FCCD\psi
  \qquad\Longrightarrow\qquad
  \mathrm{BF}_\text{res} = \Ord{\Delta x^4}
  \label{eq:bf_fccd_unification}
\end{equation}
により残差を高次へ還元する（§\ref{sec:fccd_def}）．

% -------------------------------------------------------------------------------
\subsubsection{面ジェット $\FaceJet{p}$ による BF 構築}
\label{sec:face_jet_bf_construction}
% -------------------------------------------------------------------------------

面ジェット $\FaceJet{p}=(p_f,p'_f,p''_f)$ は
P-1..P-4 を単一オブジェクトで具現する．
変密度 PPE の\textbf{面フラックス}は
\begin{equation}
  F^p_f := \betaf\,\FaceJet{p}_1 = \betaf\,p'_f,
  \label{eq:fccd_bf_face_flux}
\end{equation}
ここで $\betaf = (1/\rho)_f$（P-4 単一定義）．
PPE は FVM セル積分で
\begin{equation}
  F^p_{i+1/2} - F^p_{i-1/2}
    = \int_{x_{i-1/2}}^{x_{i+1/2}}\! S\,dx,
  \qquad S = \frac{1}{\Delta t}\,D_h^\text{bf}\,\mathbf{u}^*,
  \label{eq:fccd_bf_cell_balance}
\end{equation}
として組み立て，速度補正は
\begin{equation}
  u_f^{n+1}
    = u_f^* - \Delta t\,\betaf\,p'_f
    \quad\text{（同じ }p'_f\text{）．}
  \label{eq:fccd_bf_corrector}
\end{equation}
\textbf{同じ $p'_f$} が PPE 組立と補正で共有されるため P-2 が構造的に成立．

% -------------------------------------------------------------------------------
\subsubsection{界面ジャンプ GFM 補正（Phase 1 実用）}
\label{sec:fccd_bf_gfm_integration}
% -------------------------------------------------------------------------------

界面 $\Gamma$ を跨ぐ面 $f$ で面ジェットを
ジャンプ条件 $[p]_\Gamma=\sigma\kappa$，$[\betaf\partial_n p]_\Gamma=0$
（非相変化）で補正：
\begin{equation}
  p'^{\text{corrected}}_f
    = p'_f - \frac{\sigma\kappa}{H_f},
  \qquad p^{\text{corrected}}_f
    = p_f + \frac{\sigma\kappa}{2}\,\mathrm{sign}(\phi_f).
  \label{eq:fccd_bf_gfm_face}
\end{equation}
bulk ステンシルは不変；ジャンプベクトル $c_\Gamma$ は
$\kappa$ 更新時のみ再計算．
Phase 1 の BF 精度は 2 次，H-01 診断に十分．

\noindent\textbf{Phase 3（IIM 行補正，研究）}：
FCCD 関係 $M_f\mathbf{g} = R_f\mathbf{p} + c_\Gamma^\text{IIM}$ により
コンパクト構造を維持したまま $\Ord{h^3}$--$\Ord{h^4}$ まで昇格．

% -------------------------------------------------------------------------------
\subsubsection{静水圧テストによる BF 精度検証}
\label{sec:fccd_bf_hydrostatic}
% -------------------------------------------------------------------------------

\noindent\textbf{テスト}：重力 $\mathbf{g}=-g\hat{\mathbf{z}}$ 下の静止 2 層流体．
解析解は $p_\text{stat}(z) = -\rho(\phi(z))\,g\,z$，$\mathbf{u}=\mathbf{0}$．

\noindent\textbf{BF 残差測定}：
\begin{equation}
  \varepsilon_\text{BF}
    := \max_f\,\bigl|\,u_f^{n+1}\,\bigr|
    \quad\text{after one IPC step with }\mathbf{u}^*=\mathbf{0},\ p=p_\text{stat}.
  \label{eq:fccd_bf_hydrostatic_metric}
\end{equation}
理想離散 BF では $\varepsilon_\text{BF}=0$．実測（予測）：
\begin{center}
\begin{tabular}{@{}lll@{}}
\hline
演算子 & 非一様格子 $\rho_l/\rho_g=830$ & 収束率 \\
\hline
旧設計（節点 CCD + 事後面平均）& $\varepsilon_\text{BF}\sim 10^{-3}$ & $\Ord{\Delta x^2}$ \\
FCCD + GFM（Phase 1）& $\varepsilon_\text{BF}\sim 10^{-6}$ & $\Ord{\Delta x^2}$$^{\ddagger}$ \\
FCCD + IIM（Phase 3 目標）& $\varepsilon_\text{BF}\sim 10^{-10}$ & $\Ord{\Delta x^4}$ \\
\hline
\end{tabular}
\end{center}
\noindent$^{\ddagger}$ GFM の 2 次律速により収束率は Theorem 1（\S\ref{sec:fccd_adv_bf_theorem}）
の $\Ord{H^4}$ 上界に届かない．絶対誤差 $10^{-6}$ は一様格子 FCCD $\Ord{H^4}$
係数の小さいことから得られる実測予測値であり，漸近オーダーは GFM 依存で
$\Ord{\Delta x^2}$．IIM（Phase 3）で初めて $\Ord{\Delta x^4}$ に昇格する．

\noindent\textbf{重力 BF テスト}として明示：本表は重力平衡での BF 残差を測定する
（CSF BF の静止液滴テスト \S\ref{sec:hydrostatic_droplet} とは独立；
CSF 側は $f_\sigma$ と $G_h p_\sigma$ の相殺が主論点）．
予測値（§\ref{sec:verify_split_ppe_dc} で検証予定）．

詳細検証は \S\ref{sec:verify_split_ppe_dc}（§11 検証章）．

% -------------------------------------------------------------------------------
\subsubsection{設計戦略 A/B/C の比較}
\label{sec:fccd_bf_strategy_table}
% -------------------------------------------------------------------------------

\begin{center}
\begin{tabular}{@{}p{0.16\textwidth}p{0.42\textwidth}p{0.32\textwidth}@{}}
\hline
戦略 & 内容 & 用途 \\
\hline
\textbf{A：全面演算子} &
$G_h^\text{bf},D_h^\text{bf}$ を FCCD または低次面フラックスで構築．
BF sub-system を完全分離．
& 本稿推奨（Level 2/3）．
\\
\textbf{B：節点 CCD + BF 層} &
bulk は節点 CCD．PPE + 補正 + $f_\sigma$ にのみ $G_h^\text{bf}$ を追加．
低侵襲．
& 既存 CCD コード資産がある場合の pragmatic 選択．
\\
\textbf{C：毛管ジャンプ PPE} &
CSF 体積力を使わず，$[p]_\Gamma=\sigma\kappa$ を PPE に IIM 埋込み．
静止平衡は構造的厳密．
& 研究レベル，動圧分離が必要．
\\
\hline
\end{tabular}
\end{center}

\noindent\textbf{本稿の位置}：
Level 1（SSPRK3 検証）= 戦略 B（節点 CCD + GFM）；
Level 2（プロダクション）= 戦略 A（FCCD 面ジェット + GFM）；
Level 3（純 FCCD DNS）= 戦略 C を目標とする分相 PPE
（\S\ref{sec:split_ppe_framework}）．

% -------------------------------------------------------------------------------
\subsubsection{BF アンチパターン一覧}
\label{sec:fccd_bf_anti_patterns}
% -------------------------------------------------------------------------------

\begin{tcolorbox}[warnbox,title={避けるべき BF アンチパターン}]
\begin{itemize}[leftmargin=1.5em,noitemsep]
  \item 高次 $G_h p$，低次 $f_\sigma$（切断誤差非対称）
  \item 節点 CCD 全量 $\kappa$（界面帯ノイズ増幅）
  \item PPE $G_h$ と補正 $G_h$ が別演算子（F-2）
  \item $f_\sigma$ が節点，$G_h p$ が面（F-1）
  \item $\betaf$ 混用（PPE 調和 / 補正算術 / $f_\sigma$ 他；F-4）
  \item 界面越え CCD ステンシル無補正（F-5）
  \item BF 補正経路に Rhie--Chow（$G_h$ 二重定義）
  \item BF 経路に非対称 DCCD フィルタ（片側減衰）
\end{itemize}
\end{tcolorbox}

% ═══════════════════════════════════════════════════════════════════════════════
